\documentclass[12pt,a4paper]{article}

% 中文支持
\usepackage{ctex}

% 数学包
\usepackage{amsmath,amssymb,amsthm}
\usepackage{physics}
\usepackage{booktabs}

% 页面设置
\usepackage{geometry}
\geometry{margin=2.5cm}

% 定理环境
\theoremstyle{definition}
\newtheorem{definition}{定义}[section]
\newtheorem{theorem}{定理}[section]

\title{\textbf{声子非线性传播理论}\\[0.5em]
\large 基于非线性薛定谔方程的声子波包动力学}
\author{基于分步傅里叶方法的数值仿真}
\date{\today}

\begin{document}

\maketitle

\begin{abstract}
本文档详细推导了GHz频段声子在一维波导中的非线性传播理论。考虑9GHz泵浦(15dBm,10\%IDT效率)产生5GHz信号声子,在群速度$v_g=3800$m/s、净增益6-7dB(300$\mu$m波导)条件下的传播特性。通过非线性薛定谔方程(NLSE)描述色散、非线性、增益与损耗的竞争,并使用分步傅里叶方法(SSFM)进行数值求解。
\end{abstract}

\tableofcontents
\newpage

%=============================================================================
\section{物理模型}
%=============================================================================

\subsection{系统配置}

\begin{table}[h]
\centering
\caption{实验参数}
\begin{tabular}{lll}
\toprule
参数 & 数值 & 说明 \\
\midrule
泵浦频率 & 9 GHz & 微波泵浦 \\
泵浦功率 & 15 dBm & 约31.6 mW \\
IDT效率 & 10\% & 电-声转换效率 \\
信号声子频率 & 5 GHz & 待研究频段 \\
群速度$v_g$ & 3800 m/s & 波导中的声速 \\
波导长度$L$ & 300$\mu$m & 传播距离 \\
净增益 & 6-7 dB & 增益略大于损耗 \\
传播损耗$\alpha$ & 约4 dB/cm & 材料损耗 \\
\bottomrule
\end{tabular}
\end{table}

\subsection{非线性薛定谔方程}

在慢变包络近似(SVEA)下,声子波包$u(z,t)$的传播满足非线性薛定谔方程:

\begin{equation}
\boxed{i\frac{\partial u}{\partial z}=-\frac{\beta_2}{2}\frac{\partial^2 u}{\partial t^2}+\gamma|u|^2u+i\frac{\alpha-g}{2}u}
\label{eq:NLSE}
\end{equation}

其中各项的物理意义:
\begin{itemize}
    \item $\beta_2$:群速度色散(GVD),单位s$^2$/m
    \item $\gamma$:非线性系数,单位W$^{-1}$m$^{-1}$
    \item $\alpha$:传播损耗,单位Np/m
    \item $g$:增益系数,单位Np/m
\end{itemize}

%=============================================================================
\section{方程推导}
%=============================================================================

\subsection{从波动方程到包络方程}

\subsubsection{一维波动方程}

考虑一维弹性波在介质中的传播:
\begin{equation}
\rho\frac{\partial^2\xi}{\partial t^2}=\frac{\partial\sigma}{\partial z}
\label{eq:wave_eq}
\end{equation}
其中$\xi$为位移,$\sigma$为应力,$\rho$为密度。

\subsubsection{非线性应力-应变关系}

对于非线性介质,应力-应变关系包含高阶项:
\begin{equation}
\sigma=C_1\epsilon+C_2\epsilon^2+C_3\epsilon^3+\cdots
\label{eq:nonlinear_stress}
\end{equation}
其中$\epsilon=\partial\xi/\partial z$为应变,$C_1$为线性弹性模量,$C_2,C_3$为非线性系数。

\subsubsection{慢变包络近似}

设声子场为载波与慢变包络的乘积:
\begin{equation}
\xi(z,t)=\frac{1}{2}\left[u(z,t)e^{i(k_0z-\omega_0t)}+\text{c.c.}\right]
\label{eq:envelope_ansatz}
\end{equation}

将(\ref{eq:envelope_ansatz})代入(\ref{eq:wave_eq}),并利用(\ref{eq:nonlinear_stress}),在旋转波近似下得到包络方程。

\subsection{色散关系}

线性色散关系$\omega(k)$在载波频率附近展开:
\begin{equation}
\omega(k)=\omega_0+v_g(k-k_0)+\frac{1}{2}\frac{d^2\omega}{dk^2}(k-k_0)^2+\cdots
\label{eq:dispersion_expansion}
\end{equation}

其中:
\begin{itemize}
    \item $v_g=d\omega/dk$:群速度
    \item $\beta_2=d^2k/d\omega^2=-1/v_g^2\cdot dv_g/d\omega$:群速度色散
\end{itemize}

\subsection{非线性系数的来源}

\subsubsection{三阶非线性(Kerr效应)}

声子Kerr效应来自三阶非线性弹性:
\begin{equation}
\gamma=\frac{\omega_0^2C_3}{2\rho v_g^3A_{\text{eff}}}
\label{eq:nonlinear_coefficient}
\end{equation}
其中$A_{\text{eff}}$为有效模式面积。

%=============================================================================
\section{孤子解}
%=============================================================================

\subsection{基态孤子}

在无损、无色散情况下($\alpha=g=0,\beta_2<0$),NLSE有解析孤子解:
\begin{equation}
\boxed{u(z,t)=\sqrt{P_0}\text{sech}\left(\frac{t}{\tau_0}\right)e^{i\gamma P_0z/2}}
\end{equation}
其中:
\begin{itemize}
    \item $P_0$:峰值功率
    \item $\tau_0$:脉冲宽度(特征时间)
    \item 满足孤子条件:$N=1$,即$\gamma P_0\tau_0^2/|\beta_2|=1$
\end{itemize}

\subsection{孤子条件}

基态孤子($N=1$)要求:
\begin{equation}
\boxed{\frac{\gamma P_0\tau_0^2}{|\beta_2|}=1}
\label{eq:soliton_condition}
\end{equation}
或等价地:
\begin{equation}
P_0=\frac{|\beta_2|}{\gamma\tau_0^2}
\end{equation}

\subsection{有增益/损耗时的孤子}

考虑增益($g>\alpha$),孤子振幅按指数增长:
\begin{equation}
u(z,t)=\sqrt{P_0}e^{g_{\text{eff}}z/2}\text{sech}\left(\frac{t}{\tau_0}\right)e^{i\phi(z)}
\end{equation}
其中$g_{\text{eff}}=g-\alpha$为净增益系数。

%=============================================================================
\section{分步傅里叶方法}
%=============================================================================

\subsection{算法原理}

NLSE包含线性和非线性算符:
\begin{equation}
\frac{\partial u}{\partial z}=(\hat{L}+\hat{N})u
\end{equation}
其中:
\begin{align}
\hat{L}&=-\frac{i\beta_2}{2}\frac{\partial^2}{\partial t^2}-\frac{\alpha-g}{2}\\
\hat{N}&=i\gamma|u|^2
\end{align}

在短距离$\Delta z$内,近似:
\begin{equation}
u(z+\Delta z,t)\approx e^{\Delta z\hat{L}}e^{\Delta z\hat{N}}u(z,t)
\end{equation}

\subsection{对称分步方法}

为提高精度,使用对称分步:
\begin{equation}
u(z+\Delta z)=e^{\Delta z\hat{L}/2}e^{\Delta z\hat{N}}e^{\Delta z\hat{L}/2}u(z)
\end{equation}

\subsection{计算步骤}

\begin{enumerate}
    \item \textbf{非线性步(半步)}:时域计算$u_{1/2}=u(z,t)e^{(i\gamma|u|^2+g_{\text{eff}}/2)\Delta z/2}$
    \item \textbf{色散步(全步)}:傅里叶域计算$\tilde{u}_{3/4}=\tilde{u}_{1/2}e^{-i\beta_2\omega^2\Delta z/2}$
    \item \textbf{非线性步(半步)}:时域计算$u(z+\Delta z,t)=u_{3/4}e^{(i\gamma|u|^2+g_{\text{eff}}/2)\Delta z/2}$
\end{enumerate}

%=============================================================================
\section{Python实现}
%=============================================================================

\subsection{核心代码}

\begin{verbatim}
def ssfm_propagate(u0, z, t, params):
    u = np.zeros((len(z), len(t)), dtype=complex)
    u[0, :] = u0
    u_current = u0.copy()
    
    for n in range(1, len(z)):
        # Step 1: Nonlinear (half step)
        u_half = u_current * np.exp(
            (params.g_eff/2 + 1j*params.GAMMA_NL*np.abs(u_current)**2) 
            * dz/2)
        
        # Step 2: Dispersion (full step in frequency domain)
        u_fft = np.fft.fft(u_half)
        u_fft *= np.exp(-1j*params.BETA_2/2*omega**2*dz)
        u_disp = np.fft.ifft(u_fft)
        
        # Step 3: Nonlinear (half step)
        u_current = u_disp * np.exp(
            (params.g_eff/2 + 1j*params.GAMMA_NL*np.abs(u_disp)**2) 
            * dz/2)
        u[n, :] = u_current
    
    return u
\end{verbatim}

%=============================================================================
\section{仿真结果}
%=============================================================================

\subsection{参数设置}

\begin{itemize}
    \item 信号频率:5 GHz
    \item 群速度:$v_g=3800$ m/s
    \item 波导长度:$L=300\mu$m
    \item 净增益:6.5 dB
    \item GVD:$\beta_2=10^{-20}$ s$^2$/m
    \item 非线性系数:$\gamma=10^{-3}$ W$^{-1}$m$^{-1}$
\end{itemize}

\subsection{基态孤子传播}

当初始脉冲满足孤子条件时,脉冲在传播过程中保持形状不变(仅幅度受增益影响增长)。

\subsection{高阶孤子}

当$N>1$时,脉冲在传播过程中周期性变化(呼吸子)。

%=============================================================================
\section{结论}
%=============================================================================

本文档建立了GHz频段声子非线性传播的完整理论框架,主要结论:
\begin{enumerate}
    \item NLSE准确描述了色散、非线性、增益/损耗的竞争
    \item 基态孤子在净增益介质中保持形状并指数增长
    \item SSFM是求解NLSE的高效数值方法
\end{enumerate}

\end{document}
